

% \begin{figure}[t]
% 	%\vspace{-2pt}
% 	\begin{center}
% 		\includegraphics[width=0.5\columnwidth]{figures/ceil-framework.png}
% 	\end{center}
% 	%\vspace{-2pt}
% 	\caption{CEIL in the \textit{online} IL setting, where ${f}_{\bar{\phi}}$ is an exponential moving average (EMA) of $f_\phi$, and we use the expert matching loss to update both the embedding function $f_\phi$ and the optimal $\bz^*$. }
% 	\label{fig:test} 
% 	%\vspace{-2pt}
% \end{figure}


\section{Method}

In this section, we will formulate IL as a bi-level optimization problem, which will allow us to derive our method, contextual imitation learning (CEIL). Instead of attempting to train the learning policy $\pi_\theta(\ba|\bs)$ with the plain expert matching objective (Equation~\ref{eq:min-il}), our approach introduces an additional contextual variable $\bz$ for a contextual IL policy $\pi_\theta(\ba|\bs, \cdot)$. 
The main idea of CEIL is to learn a contextual policy $\pi_\theta(\ba|\bs, \bz)$ and an optimal contextual variable $\bz^*$ such that the given expert data (demonstrations in LfD or observations in LfO) can be recovered by the learned $\bz^*$-conditioned policy $\pi_\theta(\ba|\bs, \bz^*)$. 
We begin by describing the overall framework of CEIL in Section~\ref{sec:ceil-method}, and make a connection between CEIL and the plain expert matching objective in Section~\ref{sec:mutual-info-reg}, which leads to a practical implementation under various IL settings in Section~\ref{sec:practical-implementation}. 
%In Section 4.3, we discuss the application of CEIL on offline IL setting, LfD setting, and cross-domain IL setting. 

%\textbf{Notation.} In Section~\ref{sec:imitation-learning}, we use the non-tilde $\btau$ and the tilde $\tbtau$ to distinguish between the trajectories of learning policy $\pi_\theta$ and the unknown expert policy $\pi_E$. For brevity, in this section, we 

%we introduce an additional tilde $\tbtau$ for represent the expert demonstrations/observations, 

%(For brevity, we overload the notation $\textcolor{red}{f}$ to represent a \textit{pre-defined} ``hindsight" relabeling function, which extract different semantic information compared to the extraction function $\textcolor{red}{f}$ in the LfO objective.)  

\subsection{Contextual Imitation Learning (CEIL)}
\label{sec:ceil-method}

Motivated by the hindsight information matching in online/offline RL (Section~\ref{sec:hindsight-infomation-matching}), we propose to learn a hindsight embedding function ${f_\phi}$, which encodes trajectory $\btau$ (with window size $T$) into a latent variable $\bz \in \mathcal{Z}$, $|\mathcal{Z}| \ll T*|\mathcal{S}|$. 
Formally, we learn the embedding function ${f_\phi}$ and a corresponding contextual policy $\pi_\theta(\ba|\bs,\bz)$ by minimizing the trajectory self-consistency loss: 
\begin{align}
	\pi_\theta, f_\phi &= \min_{\pi_\theta, f_\phi} \ - \mathbb{E}_{\mathcal{D}(\btau)} \lbr \log \pi_\theta(\btau|{f_\phi}(\btau)) \rbr = \min_{\pi_\theta, f_\phi} \ - \mathbb{E}_{\btau\sim\mathcal{D}(\btau)}\mathbb{E}_{(\bs,\ba)\sim\btau} \lbr \log \pi_\theta(\ba|\bs, {f_\phi}(\btau)) \rbr, \label{eq:traj-self-consistent} 
\end{align}
where in the online setting, we sample trajectory $\btau$ from buffer $\mathcal{D}(\btau)$, known as the experience replay buffer in online RL; in the offline setting, we sample trajectory $\btau$ directly from the given offline data. 

If we can ensure that the learned contextual policy $\pi_\theta$ and the embedding function $f_\phi$ are accurate on the empirical data $\mathcal{D}(\btau)$, then we can convert the IL policy optimization objective (in Equation~\ref{eq:min-il}) into a bi-level expert matching objective: 
\begin{align}
	&\min_{\bz^*} \ D(\pi_\theta(\btau|\bz^*), \pi_E(\btau)), \label{eq:ceil-obj} \\ 
	&\quad \text{s.t. } \pi_\theta, f_\phi = \min_{\pi_\theta, f_\phi}  - \mathbb{E}_{\mathcal{D}(\btau)}  \lbr \log \pi_\theta(\btau|{f_\phi}(\btau))\rbr - \mathcal{R}(f_\phi), \ \  \text{and} \ \  \bz^* \in {f_\phi} \circ \texttt{supp}(\mathcal{D}), \label{eq:ceil-obj-supp}
\end{align}
where $\mathcal{R}(f_\phi)$ is an added regularization over the embedding function (we will elaborate on it later), and $\texttt{supp}(\mathcal{D})$ denotes the support of the trajectory distribution $\{ \btau | \mathcal{D}(\btau)>0 \}$. Here ${f_\phi}$ is employed to map the trajectory space to the latent variable space ($\mathcal{Z}$). 
Intuitively, by optimizing Equation~\ref{eq:ceil-obj}, we expect the induced trajectory distribution of the learned $\pi_\theta(\ba|\bs, \bz^*)$ will match that of the expert. 
However, 
in the offline IL setting, the contextual policy can not interact with the environment. 
If we directly optimize the expert matching objective (Equation~\ref{eq:ceil-obj}), such an objective can easily exploit generalization errors in the contextual policy model to infer a mistakenly overestimated $\bz^*$ that achieves low expert-matching loss but does not preserve the trajectory self-consistency (Equation~\ref{eq:traj-self-consistent}). Therefore, we formalize CEIL into a bi-level optimization problem, where, in Equation~\ref{eq:ceil-obj-supp}, we explicitly constrain the inferred $\bz^*$ lies in the (${f_\phi}$-mapped) support of the training trajectory distribution. 

At a high level, CEIL decouples the learning policy into two parts: an expressive contextual policy $\pi_\theta(\ba|\bs, \cdot)$ and an optimal contextual variable $\bz^*$. 
By comparing CEIL with the plain expert matching objective, $\min_{\pi_\theta} D(\pi_\theta(\btau), \pi_E(\btau))$, in Equation~\ref{eq:min-il}, we highlight two merits: 1) CEIL's expert matching loss (Equation~\ref{eq:ceil-obj}) does not account for updating $\pi_\theta$ and is only incentivized to update the low-dimensional latent variable $\bz^*$, which enjoys efficient parameter learning similar to the prompt tuning in large language models~\citep{zhou2022learning}, and 2) we learn $\pi_\theta$ by simply performing supervised regression (Equation~\ref{eq:ceil-obj-supp}), which is more stable compared to vanilla inverse-RL/adversarial-IL methods.






\subsection{Connection to the Plain Expert Matching Objective}
\label{sec:mutual-info-reg}

To gain more insight into Equation~\ref{eq:ceil-obj} that captures the quality of IL (the degree of similarity to the expert data), we define $D(\cdot, \cdot)$ as the sum of reverse KL and forward KL divergence\footnote{$D_\text{KL}(p\Vert q):=\mathbb{E}_{p(\bx)}\left[ \log \frac{p(\bx)}{q(\bx)} \right]$ denotes the (forward) KL divergences. It is well known that reverse KL ensures that the learned distribution is mode-seeking and forward KL exhibits a mode-covering behavior~\citep{ke2020imitation}. For analysis purposes, here we define $D(\cdot, \cdot)$ as the sum of reverse KL and forward KL, and set the weights of both reverse KL and forward KL to 1.}, \ie $D(q, p) = D_\text{KL}(q\Vert p) + D_\text{KL}(p\Vert q)$, and derive an alternative form for Equation~\ref{eq:ceil-obj}: 
\begin{align}
	&\ \quad \arg \min_{\bz^*} \ D(\pi_\theta(\btau|\bz^*), \pi_E(\btau)) 
	= \arg \max_{\bz^*} \ \underbrace{\mathcal{I}(\bz^*; \btau_E) - \mathcal{I}(\bz^*; \btau_\theta)}_{\mathcal{J}_\text{MI}}
	- \underbrace{D_\text{KL}(\pi_\theta(\btau), \pi_E(\btau))}_{\mathcal{J}_{D}}, \label{eq:ceil-mi-reg}
\end{align}
where $\mathcal{I}(\bx;\by)$ denotes the mutual information (MI) between $\bx$ and $\by$, which measures the predictive power of $\by$ on $\bx$ (or vice-versa), 
the latent variables are defined as $\btau_E:=\btau \sim \pi_E(\btau)$, $\btau_\theta:=\btau \sim p(\bz^*)\pi_\theta(\btau|\bz^*)$, and $\pi_\theta(\btau)=\mathbb{E}_{\bz^*}\left[ \pi_\theta(\btau|\bz^*) \right]$. 

Intuitively, the second term $\mathcal{J}_{D}$ on RHS of Equation~\ref{eq:ceil-mi-reg} is similar to the plain expert matching objective in Equation~\ref{eq:min-il}, except that here we optimize a latent variable $\bz^*$ over this objective. Regarding the MI terms $\mathcal{J}_\text{MI}$, we can interpret the maximization over $\mathcal{J}_\text{MI}$ as an implicit policy improvement, which incentivizes the optimal latent variable $\bz^*$ for having high predictive power of the expert data $\btau_E$ and having low predictive power of the non-expert data $\btau_\theta$. 

Further, we can rewrite the MI term ($\mathcal{J}_\text{MI}$ in Equation~\ref{eq:ceil-mi-reg}) in terms of the learned embedding function $f_\phi$, yielding an approximate embedding inference objective $\mathcal{J}_{\text{MI}(f_\phi)}$: % analogous to contrastive learning:
\begin{align}
	\mathcal{J}_\text{MI} 
	&= \mathbb{E}_{\pi_E(\bz^*, \btau_E)}\log p(\bz^*|\btau_E) - \mathbb{E}_{\pi_\theta(\bz^*, \btau_\theta)}\log p(\bz^*|\btau_\theta)  \nonumber \\
	%	&= \mathbb{E}_{\delta(\bz^*)\pi_E(\btau_E)\pi_\theta(\btau_E)} \lbr \log p(\bz^*|\bz_E) + \log p(\bz_E|\btau_E)\right. \nonumber\\
	%	&\ \quad \quad \quad \quad \quad \quad \quad \quad \ \  \left. -  \log p(\bz^*|\bz_\theta) - \log p(\bz_\theta|\btau_\theta)\rbr \nonumber\\
	%	&\approx \mathbb{E}_{\delta(\bz^*)\pi_E(\btau_E)\pi_\theta(\btau_E)} \lbr - \Vert\bz^*-\bz_E\Vert^2 + \Vert\bz^*-\bz_\theta\Vert^2 \rbr \nonumber,
	&\approx \mathbb{E}_{p(\bz^*)\pi_E(\btau_E)\pi_\theta(\btau_\theta|\bz^*)} \lbr - \Vert\bz^*-f_\phi(\btau_E)\Vert^2 + \Vert\bz^*-f_\phi(\btau_\theta)\Vert^2 \rbr \nonumber \triangleq \mathcal{J}_{\text{MI}(f_\phi)}, 
\end{align} %\bz^*, \btau_E, \btau_\theta
where we approximate the logarithmic predictive power of $\bz^*$ on $\btau$ with $- \Vert\bz^*-f_\phi(\btau)\Vert^2$, by taking advantage of the learned embedding function $f_\phi$ in Equation~\ref{eq:ceil-obj-supp}. 




By maximizing $\mathcal{J}_{\text{MI}(f_\phi)}$, %in the latent space,  
the learned optimal $\bz^*$ will be induced to converge towards the embeddings of expert data and avoid trivial solutions (as shown in Figure~\ref{fig:ceil-toy}). 
Intuitively, 
$\mathcal{J}_{\text{MI}(f_\phi)}$ can also be thought of as an instantiation of contrastive loss, which manifests two facets we consider~significant~in~IL: 

\begin{wrapfigure}{r}{0.525\textwidth}
	\vspace{-1pt}
	\small 
	\begin{center}
		\includegraphics[scale=0.345]{figures/anchor/anchor-hopper.pdf} \ 
		\includegraphics[scale=0.345]{figures/anchor/anchor-ant.pdf}
		\includegraphics[scale=0.345]{figures/anchor/anchor_legend.pdf}
	\end{center}
	\vspace{-5pt}
	\caption{{During learning, the distance between $\bz^*$ and $f_\phi(\btau_E)$ decreases rapidly (green lines). Meanwhile, as policy $\pi_\theta(\cdot|\cdot, \bz^*)$ gets better (blue lines), $f_\phi(\btau_\theta)$ gradually approaches $\bz^*$ (red lines). }
	}
	\vspace{-6pt}
	\label{fig:ceil-toy}
\end{wrapfigure}
1) the "anchor" variable\footnote{The triplet contrastive loss enforces the distance between the anchor and the positive to be smaller than that between the anchor and the negative. Thus, we can view $\bz^*$ in $\mathcal{J}_{\text{MI}(f_\phi)}$ as an instance of the anchor.} $\bz^*$ is unknown and must be estimated, and 2) it is necessary to ensure that the estimated $\bz^*$ lies in the support set of training distribution, as specified by the support constraints in Equation~\ref{eq:ceil-obj-supp}. 




In summary, by comparing $\mathcal{J}_{\text{MI}(f_\phi)}$ and $\mathcal{J}_{D}$, we can observe that $\mathcal{J}_{\text{MI}(f_\phi)}$ actually encourages expert matching in the embedding space, while $\mathcal{J}_{D}$ encourages expert matching in the original trajectory space. In the next section, we will see that such an embedding-level expert matching objective naturally lends itself to cross-domain IL settings. 


\begin{algorithm}[t]
	\caption{Training CEIL:  \textcolor{red}{Online} \textbf{\textit{or}} \textcolor{blue}{Offline} IL Setting}
	\label{alg:ceil}
	\textbf{Require:} Expert demonstrations $\{\pi_E(\btau)\}$, \textcolor{red}{empty buffer $\mathcal{D}$ for {online} IL} \textbf{\textit{or}} \textcolor{blue}{reward-free offline data $\mathcal{D}$ for {offline} IL}, training iteration $K$, and batch size $N$. 
	
	\begin{algorithmic}[1]
		\STATE Initialize contextual policy $\pi_\theta(\ba|\bs,\cdot)$, embedding function $f_\phi(\bz|\btau)$, and latent variable $\bz^*$. 
		\FOR{$k=1,\cdots, K$}
		\STATE \textcolor{red}{(Online only) Run policy $\pi_\theta(\ba|\bs,\bz^*)$ in environment and store experience into buffer $\mathcal{D}$.} % $\{\pi_\theta(\btau|\bz^*)\}$ 
		\STATE Sample a batch of data $\{ \btau \}_1^n$ from \textcolor{red}{$\mathcal{D}$ for online IL} \textbf{\textit{or}} \textcolor{blue}{$\mathcal{D}$ for offline IL}. % \\ \textcolor{gray}{/* Lines 5 and 6 are optimized over sampled $\{ \btau \}_1^n$ */}
		\STATE Learn $\pi_\theta$ and $f_\phi$ over sampled $\{ \btau \}_1^n$ using the trajectory self-consistency loss. 
		\STATE Update $\bz^*$ and $f_\phi$ over sampled $\{ \btau \}_1^n$ by maximizing $\mathcal{J}_{\text{MI}(f_\phi)} - \alpha \mathcal{J}_{D}$.
		\STATE \textcolor{blue}{(Offline only) Update $\bz^*$ by minimizing  $\mathcal{R}(\bz^*)$.} \textcolor{gray}{\ \ \# eliminating the offline OOD issues.} %constrain $\bz^*$ to stay in the offline support set.} % Equation~\ref{eq:ceil-stop-grad}. 
		\ENDFOR
		%		\STATE Train the diffusion prior $\pi_\psi$ with Equation~\ref{eq:didi-prior}.
		%		\STATE Train the return network $R(\btau_t)$ of trajectory in $\pi_\mathcal{D}(\btau_t)$. 
		%		\WHILE{not converged}
		%		\STATE Sample a skill, $\bz \sim p(\bz)$, and a state $\bs_t \sim \pi_\mathcal{D}(\btau_t)$.
		%		\STATE Learn $q_\phi(\bz|\textcolor{minecolortau}{\btau^n_t})$ and $\pi_\theta(\textcolor{minecolortau}{\btau^n_t}|\bs_t, \bz)$ with $\mathcal{J}_\text{DIDI}$. 
		%		\ENDWHILE
	\end{algorithmic}
	\textbf{Return:} the learned contextual policy $\pi_\theta(\ba|\bs,\cdot)$ and the optimal latent variable $\bz^*$.
\end{algorithm}
\subsection{Practical Implementation}
\label{sec:practical-implementation}

In this section, we describe how we can convert the bi-level IL problem above (Equations~\ref{eq:ceil-obj} and \ref{eq:ceil-obj-supp}) into a feasible online/offline IL objective and discuss some practical implementation details in LfO, offline IL, cross-domain IL, and one-shot IL settings ({see more details in Appendix~\ref{app:ceil_implementation_details} and our code will be released at \url{https://github.com/wechto/GeneralizedCEIL}.}). 


As shown in Algorithm~\ref{alg:ceil} (best viewed in colors), CEIL alternates between solving the bi-level problem with respect to the support constraint (Line 3 for \textcolor{red}{online} IL or Line 7 for \textcolor{blue}{offline} IL), the trajectory self-consistency loss (Line 5), and the optimal embedding inference (Line 6). 

To satisfy the support constraint in Equation~\ref{eq:ceil-obj-supp}, for \textcolor{red}{online} IL (Line 3), we directly roll out the $z^*$-conditioned policy $\pi_\theta(\ba|\bs,\bz^*)$ in the environment; for \textcolor{blue}{offline} IL (Line 7), we minimize a simple regularization over $\bz^*$, bearing a close resemblance to the one used in TD3+BC~\citep{fujimoto2021minimalist}:  
% constrain the learned embedding $\bz^*$ to stay in the support set:
% \begin{align}
% 	\min_{\bz^*} \ \mathbb{E}_{p(\bz^*)\pi_\theta(\btau|\bz^*)} \left[ \Vert \bz^* - f_{\bar{\phi}}(\btau)\Vert^2 \right], %\label{eq:ceil-stop-grad}
% \end{align}
\begin{align}\label{eq:ceil-stop-grad}
    \mathcal{R}(\bz^*) = \min \left( \Vert \bz^* - f_{\bar{\phi}}(\btau_E)\Vert^2 ,
     \Vert \bz^* - f_{\bar{\phi}}(\btau_\mathcal{D})\Vert^2 
 \right), \ \ \btau_E:=\btau \sim \pi_E(\btau), \  \btau_\mathcal{D}:=\btau \sim \mathcal{D}(\btau), 
\end{align}
where we apply a stop-gradient operation to $f_{\bar{\phi}}$. 
To ensure the optimal embedding inference ($\max_{\bz^*} \mathcal{J}_{\text{MI}(f_\phi)} - \mathcal{J}_{D}$) retaining the flexibility of seeking $\bz^*$ across different instances of $f_\phi$, 
we jointly update the optimal embedding $\bz^*$ and the embedding function $f_\phi$ with
\begin{align}\label{eq:ceil-joint-opt}
 \max_{\bz^*, f_\phi} \  \mathcal{J}_{\text{MI}(f_\phi)} - \alpha \mathcal{J}_{D}, 
\end{align}
where we use $\alpha$ to control the weight on $\mathcal{J}_{D}$. 
% \begin{equation} \label{eq:ceil-joint-opt}
% \left\{
% \begin{aligned}
% &\max_{\bz^*, f_\phi} \  \mathcal{J}_{\text{MI}(f_\phi)} - \alpha_1 \mathcal{J}_{D}, \text{\quad \quad \quad \quad \quad \quad if \textcolor{red}{online} IL}, \\
% &\max_{\bz^*, f_\phi} \  \mathcal{J}_{\text{MI}(f_\phi)} - \alpha_1 \mathcal{J}_{D} - \alpha_2 \mathcal{R}(\bz^*),  \text{\quad if \textcolor{blue}{offline} IL}, 
% \end{aligned}
% \right. 
% \end{equation}
% where we use $\alpha_1$ and $\alpha_2$ to control the weights on $\mathcal{J}_{D}$ and $\mathcal{R}(\bz^*)$ respectively. 


\textbf{LfO.} In the LfO setting, as expert actions are missing, we apply our expert matching objective only over the observations. % The state space across all tasks is of the same dimensionality,. 
Note that even though expert data contains no actions in LfO, we can still leverage a large number of suboptimal actions presented in online/offline $\mathcal{D(\btau)}$. 
Thus, we can learn the contextual policy $\pi_\theta(\ba|\bs,\bz)$ using the buffer data in online IL or the offline data in offline IL, much owing to the fact that we do not directly use the plain expert matching objective to update $\pi_\theta$. 

\textbf{Cross-domain IL.} Cross-domain IL considers the case in which the expert’s and learning agent’s MDPs are different. %, such as when the learning and expert policies are of different embodiments or are in different environments (dynamics). 
Due to the domain shift, the plain idea of $\min \mathcal{J}_D$ may not be a sufficient proxy for the expert matching objective, as there may never exist a trajectory (in the learning MDP) that matches the given expert data. 
Thus, we can set (the weight of $\mathcal{J}_D$) $\alpha$ to $0$. 

Further, to make embedding function $f_\phi$ useful for guiding the expert matching in latent space (\ie $\max \mathcal{J}_{\text{MI}(f_\phi)}$), %can alternatively regularize $f_\phi$ to capture environment/domain-agnostic while task-relevant embeddings. following~\citet{franzmeyer2022learn}, Inspired by the intuition of mutual information interpretation of $\mathcal{J}_{\text{MI}}$, 
we encourage $f_\phi$ to capture the task-relevant embeddings and ignore the domain-specific factors. 
To do so, we generate a set of pseudo-random transitions $\{\btau_{E'}\}$ by independently sampling trajectories from expert data $\{\pi_E(\btau_{E})\}$ and adding random noise over these sampled trajectories, \ie $\btau_{E'} = \btau_{E} + \text{noise}$. Then, we couple each trajectory $\btau$ in $\{\btau_{E}\}\cup\{\btau_{E'}\}$ with a label $\mathbf{n} \in \{\mathbf{0}, \mathbf{1}\}$, indicating whether it is noised, and then generate a new set of $\{(\btau, \mathbf{n})\}$, where $\btau \in \{\btau_{E}\}\cup\{\btau_{E'}\}$ and $\mathbf{n} \in \{\mathbf{0}, \mathbf{1}\}$. Thus, we can set the regularization $\mathcal{R}(f_\phi)$ in Equation~\ref{eq:ceil-obj-supp} to be:
\begin{align}
	\mathcal{R}(f_\phi)  = \mathcal{I}(f_\phi(\btau); \mathbf{n}).
\end{align} 
Intuitively, maximizing $\mathcal{R}(f_\phi)$ encourages embeddings to be domain-agnostic and task-relevant: $f_\phi(\btau_E)$ has high predictive power over expert data ($\mathbf{n}=0$) and low that over noised data ($\mathbf{n}=\mathbf{1}$). 


\textbf{One-shot IL.} Benefiting from the separate design of  the contextual policy learning and the optimal embedding inference, CEIL also enjoys another advantage --- one-shot generalization to new IL tasks. 
For new IL tasks, given the corresponding expert data ${\btau}_\text{new}$, we can use the learned embedding function $f_\phi$ to generate a corresponding latent embedding ${\mathbf{{z}}}_\text{new}$. When conditioning on such an embedding, we can directly roll out $\pi_\theta(\ba|\bs, {\mathbf{{z}}}_\text{new})$ to recover the one-shot expert behavior. 

% in the environment. 



